%============================================================
% CH1 · LESSON 5 — MODELLING: SEEING THE GRAPH
%============================================================
\subsection{Modelling: Seeing the Graph}
\label{ssec:ch1-modelling}

Few contest problems say ``here is a graph''. Most describe a board, a number that can be changed, a word that can be edited, a lock, a game --- and the graph is hidden inside. Finding it is the most important step of every solution in this book, which is why each solved problem has a \emph{Graph Model} box. This lesson teaches the habit: \textbf{decide what a vertex is, decide what an edge is, estimate how many there are, and recognise the shape.} Then a standard algorithm --- often the DFS or BFS of Lesson~\ref{ssec:ch1-first-traversal} --- does the rest.

\begin{definitionbox}[Implicit graph]
A graph is \textbf{implicit} when it is not given as a list of edges but by a \emph{rule}: the vertices are all objects of some kind (cells, numbers, strings, configurations), and a function \texttt{neighbours(x)} computes the edges out of $x$ on demand. The graph is never stored; a traversal only needs \texttt{visited} (or \texttt{dist}) for the vertices it actually reaches.
\end{definitionbox}

%------------------------------------------------------------
\subsubsection*{Grids: cells are vertices, shared sides are edges}

A grid of $R\times C$ cells is the most common implicit graph. Each free cell is a vertex; with \textbf{4-neighbourhood} two cells are adjacent if they share a side, with \textbf{8-neighbourhood} also if they share a corner (king moves). The statement decides which: ``up, down, left or right'' means 4, ``including diagonals'' or ``touching at a corner'' means 8.

\begin{diagrambox}
\begin{tikzpicture}[x=0.8cm, y=0.8cm]
    \begin{scope}
        \foreach \r in {0,1,2} \foreach \c in {0,1,2} \node[gcell] at (\c,-\r) {};
        \foreach \c/\r in {1/0, 1/2, 0/1, 2/1} \node[gcell, fill=hitfill] at (\c,-\r) {};
        \node[gcell, fill=currentfill] at (1,-1) {};
        \foreach \dx/\dy in {0/1, 0/-1, -1/0, 1/0} \draw[dep, line width=1.2pt] (1+0.2*\dx,-1+0.2*\dy) -- (1+0.75*\dx,-1+0.75*\dy);
        \node[font=\small\bfseries] at (1,-3) {4 neighbours};
        \node[font=\scriptsize\ttfamily] at (1,-3.55) {dr = \{-1,1,0,0\}};
        \node[font=\scriptsize\ttfamily] at (1,-3.95) {dc = \{0,0,-1,1\}};
    \end{scope}
    \begin{scope}[xshift=3.8cm]
        \foreach \r in {0,1,2} \foreach \c in {0,1,2} \node[gcell, fill=hitfill] at (\c,-\r) {};
        \node[gcell, fill=currentfill] at (1,-1) {};
        \foreach \dx/\dy in {0/1, 0/-1, -1/0, 1/0, 1/1, 1/-1, -1/1, -1/-1} \draw[dep, line width=1.2pt] (1+0.2*\dx,-1+0.2*\dy) -- (1+0.75*\dx,-1+0.75*\dy);
        \node[font=\small\bfseries] at (1,-3) {8 neighbours};
        \node[font=\scriptsize\ttfamily] at (1,-3.55) {dr, dc in -1..1};
        \node[font=\scriptsize\ttfamily] at (1,-3.95) {skip dr = dc = 0};
    \end{scope}
    \begin{scope}[xshift=8cm, yshift=0.2cm]
        \foreach \r in {0,1,2} \foreach \c in {0,1,2} \draw[ge] (\c*0.9,-\r*0.9) -- (\c*0.9+0.9,-\r*0.9);
        \foreach \r in {0,1} \foreach \c in {0,1,2,3} \draw[ge] (\c*0.9,-\r*0.9) -- (\c*0.9,-\r*0.9-0.9);
        \foreach \r in {0,1,2} \foreach \c in {0,1,2,3} \fill[primaryblue] (\c*0.9,-\r*0.9) circle (2.6pt);
        \node[font=\small\bfseries] at (1.35,-2.4) {$3\times4$ grid as a graph};
        \node[font=\scriptsize] at (1.35,-2.95) {$|V|=12$, $|E|=17$ (4 nbrs)};
        \node[font=\scriptsize] at (1.35,-3.35) {$|E|=29$ with 8 neighbours};
    \end{scope}
\end{tikzpicture}
\figcap{The two neighbourhoods and the graph of a $3\times 4$ grid with no walls: $3\cdot3$ horizontal $+\,2\cdot4$ vertical $=17$ edges; adding the $2\cdot(2\cdot3)=12$ diagonals gives 29.}
\end{diagrambox}

In general an $R\times C$ grid has $RC$ vertices and fewer than $2RC$ edges (4-neighbourhood) --- always sparse. A $1000\times1000$ grid has $10^6$ vertices and $1\,998\,000$ edges: one BFS is a few million steps, but a DFS on it can recurse $10^6$ levels deep (Lesson~\ref{ssec:ch1-first-traversal}). Pattern~3 is devoted to grids.

%------------------------------------------------------------
\subsubsection*{Numbers: values are vertices, operations are edges}

``You start with $a$. In one move you may add 1 or multiply by 3. What is the fewest number of moves to reach $b$?'' The vertices are the integers; each operation is a directed edge $x\to x+1$ and $x\to 3x$. Fewest moves $=$ shortest path in an unweighted graph $=$ BFS.

\begin{diagrambox}
\begin{tikzpicture}[x=1.05cm, y=1cm]
    \foreach \i in {1,...,12} {
        \ifnum\i=1 \node[gvsrc, minimum size=6.5mm] (n\i) at (\i,0) {\i};
        \else\ifnum\i=10 \node[gvdst, minimum size=6.5mm] (n\i) at (\i,0) {\i};
        \else \node[gv, minimum size=6.5mm] (n\i) at (\i,0) {\i}; \fi\fi
    }
    \foreach \i [evaluate=\i as \j using int(\i+1)] in {1,...,11} \draw[gde] (n\i) -- (n\j);
    \draw[gde, draw=softpurple] (n2) to[bend left=45] (n6);
    \draw[gde, draw=softpurple] (n4) to[bend right=35] (n12);
    \draw[gpath, -{Stealth[length=2.6mm]}] (n1) to[bend left=50] (n3);
    \draw[gpath, -{Stealth[length=2.6mm]}] (n3) to[bend left=40] (n9);
    \draw[gpath, -{Stealth[length=2.6mm]}] (n9) -- (n10);
    \foreach \i/\d in {1/0, 2/1, 3/1, 9/2, 10/3} \node[gdist, below=1pt of n\i] {\d};
    \node[font=\scriptsize\color{softpurple}] at (8,1.75) {$\times 3$ arcs};
    \node[font=\scriptsize\color{darkgray}] at (6.5,-0.75) {$+1$ arcs};
\end{tikzpicture}
\figcap{The first 12 vertices of the ``$+1$ / $\times3$'' graph ($\times3$ arcs curved; the ones used by the answer are orange). The shortest route from 1 to 10 is $1\to3\to9\to10$ (3 moves; BFS distances in green). Walking with $+1$ only would take 9 moves.}
\end{diagrambox}

The neighbours are computed, not stored. One extra thought is needed: \emph{which values can ever be useful?} Both operations only increase $x$, so values above $b$ can be cut off, and the graph has at most $b$ vertices and $2b$ edges.

\begin{lstlisting}[style=cpp]
#include <bits/stdc++.h>
using namespace std;

// An implicit graph on numbers: from x you may go to x+1 or to 3x.
// Fewest operations from a to b (1 <= a <= b <= 10^6). The graph is never stored:
// the neighbours of x are computed on the fly.
int main() {
    int a, b;
    if (!(cin >> a >> b)) return 0;
    const int LIM = b;                       // larger values can never come back down
    vector<int> dist(LIM + 1, -1), par(LIM + 1, -1);
    queue<int> q;
    dist[a] = 0;
    q.push(a);
    while (!q.empty()) {
        int x = q.front();
        q.pop();
        long long nb[2] = {x + 1LL, 3LL * x};           // the two edges out of x
        for (int k = 0; k < 2; k++) {
            if (nb[k] > LIM) continue;                  // outside the useful range
            int y = (int)nb[k];
            if (dist[y] == -1) {
                dist[y] = dist[x] + 1;
                par[y] = x;
                q.push(y);
            }
        }
    }
    cout << "operations: " << dist[b] << "\npath:";
    vector<int> path;
    for (int x = b; x != -1; x = par[x]) path.push_back(x);
    reverse(path.begin(), path.end());
    for (int x : path) cout << " " << x;
    cout << "\n";
    return 0;
}
\end{lstlisting}

\begin{lstlisting}[style=trace]
input : 1 31      output: operations: 5   path: 1 3 9 10 30 31
input : 2 100     output: operations: 7   path: 2 3 9 10 11 33 99 100
input : 5 5       output: operations: 0   path: 5
\end{lstlisting}

The array \texttt{par} remembers through which edge each vertex was discovered; walking it back from $b$ rebuilds a shortest path. This ``parent array'' trick works for every BFS.

\begin{warningbox}[Bound the state space before you code]
If an operation can \emph{decrease} the value or if values may temporarily exceed the target, the cut-off is less obvious. Prove a bound (for example ``it never helps to go above $2b$'') or the BFS either misses the answer or runs out of memory. Pattern~4 works through these arguments.
\end{warningbox}

%------------------------------------------------------------
\subsubsection*{Strings: words are vertices, small edits are edges}

Words of equal length are adjacent when they differ in exactly one letter. Changing one letter at a time, \textsf{cold} becomes \textsf{warm} in four steps:

\begin{diagrambox}
\begin{tikzpicture}[x=2.1cm, y=1cm,
    w/.style={rectangle, rounded corners=3pt, draw=darkgray!85, fill=white, line width=0.9pt, font=\small\sffamily\bfseries, minimum width=13mm, minimum height=7mm}]
    \node[w, draw=successgreen, fill=basefill] (a) at (0,0) {cold};
    \node[w] (b) at (1,0) {cord};
    \node[w] (c) at (2,0) {card};
    \node[w] (d) at (3,0) {ward};
    \node[w, draw=accentorange, fill=answerfill] (e) at (4,0) {warm};
    \draw[gpath] (a) -- node[above=3pt, font=\scriptsize, text=darkgray] {l$\to$r} (b);
    \draw[gpath] (b) -- node[above=3pt, font=\scriptsize, text=darkgray] {o$\to$a} (c);
    \draw[gpath] (c) -- node[above=3pt, font=\scriptsize, text=darkgray] {c$\to$w} (d);
    \draw[gpath] (d) -- node[above=3pt, font=\scriptsize, text=darkgray] {d$\to$m} (e);
    \node[w, draw=midgray, text=midgray] (f) at (0.5,-1.3) {colt};
    \node[w, draw=midgray, text=midgray] (g) at (2,-1.3) {cart};
    \node[w, draw=midgray, text=midgray] (h) at (3.5,-1.3) {word};
    \draw[gfaded] (a) -- (f); \draw[gfaded] (c) -- (g); \draw[gfaded] (d) -- (h); \draw[gfaded] (b) -- (h);
\end{tikzpicture}
\figcap{Part of the word graph. A shortest path from \textsf{cold} to \textsf{warm} has 4 edges; grey words are other neighbours that a BFS would also explore (\textsf{word} is adjacent to both \textsf{cord} and \textsf{ward}).}
\end{diagrambox}

How to generate the neighbours of a word of length $L$? Either compare it with every word of the dictionary ($\bigO{N\cdot L}$ per word, $\bigO{N^2L}$ in total), or try all $25L$ one-letter changes (100 candidates for a 4-letter word) and look each one up in a hash set. Which is faster depends on $N$ and $L$ --- exactly the kind of estimate this lesson is about.

%------------------------------------------------------------
\subsubsection*{Game and puzzle states: tuples encoded as integers}

When the ``position'' of a puzzle consists of several numbers --- two tokens on a board, the amounts of water in some jugs, a set of switched-on lamps --- the vertex is the whole \emph{tuple}. To use plain arrays for \texttt{visited} and \texttt{dist}, encode the tuple as one integer, like the digits of a number in a mixed base.

\begin{diagrambox}
\begin{tikzpicture}[x=0.55cm, y=0.55cm]
    \foreach \r in {0,...,7} \foreach \c in {0,...,7} {
        \pgfmathtruncatemacro{\parity}{mod(\r+\c,2)}
        \ifnum\parity=0 \fill[lightgray] (\c,-\r) rectangle (\c+1,-\r-1); \fi
    }
    \draw[darkgray] (0,0) rectangle (8,-8);
    \node[gvA, minimum size=4.5mm, font=\scriptsize\bfseries] at (5.5,-2.5) {A};
    \node[gvB, minimum size=4.5mm, font=\scriptsize\bfseries] at (0.5,-7.5) {B};
    \foreach \c in {0,...,7} \node[font=\tiny\color{darkgray}] at (\c+0.5,0.45) {\c};
    \foreach \r in {0,...,7} \node[font=\tiny\color{darkgray}] at (-0.45,-\r-0.5) {\r};
    \begin{scope}[xshift=5.6cm, yshift=-0.5cm]
        \node[anchor=west, font=\small] at (0,0) {state $=(r_A, c_A, r_B, c_B) = (2,5,7,0)$};
        \node[anchor=west, font=\small] at (0,-1.4) {$\mathit{id} = ((r_A\cdot8 + c_A)\cdot 8 + r_B)\cdot 8 + c_B$};
        \node[anchor=west, font=\small] at (0,-2.6) {$\phantom{\mathit{id}} = ((2\cdot8+5)\cdot8+7)\cdot8+0 = \mathbf{1400}$};
        \node[anchor=west, font=\small] at (0,-3.9) {decode: $c_B=\mathit{id}\bmod 8$, $r_B=\lfloor \mathit{id}/8\rfloor \bmod 8$, \ldots};
        \node[anchor=west, font=\small] at (0,-5.2) {number of states: $8^4 = 4096$};
        \node[anchor=west, font=\scriptsize\color{darkgray}, text width=7cm] at (0,-6.7) {one move of A or B changes the state: an edge to another id, generated on the fly by decoding, moving, re-encoding};
    \end{scope}
\end{tikzpicture}
\figcap{Two tokens on an $8\times8$ board: the vertex is the pair of positions, encoded in base 8. Decoding 1400 gives back $(2,5,7,0)$.}
\end{diagrambox}

For a set of $k$ switches use a bitmask ($2^k$ states); for small permutations use a \texttt{map} or a rank; for two jugs with capacities $A$ and $B$ use $\mathit{id}=x\cdot(B+1)+y$. The rule is the same: \textbf{the id must be unique per state and small enough to index an array.}

%------------------------------------------------------------
\subsubsection*{The modelling checklist}

\begin{checklistbox}[What is a vertex? What is an edge?]
\begin{enumerate}[leftmargin=*]
    \item \textbf{Vertex}: what does a ``situation'' consist of? Everything that influences the future must be in it (position \emph{and} keys collected \emph{and} direction faced, \ldots), and nothing else.
    \item \textbf{Edge}: what single action changes the situation? Is it reversible (undirected) or not (directed)? Does it have a cost (weighted) or does every action count as 1?
    \item \textbf{Question}: is it reachability (DFS/BFS), fewest actions (BFS), cheapest actions (Dijkstra, Ch.~7), number of groups (components, Ch.~4), a cycle, an order (Ch.~5)?
    \item \textbf{Size}: $|V|$ and $|E|$ --- do \texttt{visited}/\texttt{dist} arrays of size $|V|$ fit in memory, and does $\bigO{|V|+|E|}$ fit in the time limit?
    \item \textbf{Explicit or implicit}: store edges only when they are given; otherwise generate neighbours on the fly.
\end{enumerate}
\end{checklistbox}

%------------------------------------------------------------
\subsubsection*{Estimating $|V|$ and $|E|$ from the constraints}

A C++ program does roughly $10^8$ simple operations per second. Before coding a traversal, multiply out the size of the state space:

\begin{table}[H]
\centering\small
\begin{tabularx}{\linewidth}{@{}X l r r l@{}}
\toprule
\textbf{Model} & \textbf{vertex} & $|V|$ & $|E|$ (arcs) & \textbf{verdict} \\
\midrule
$1000\times1000$ grid, 4 neighbours & cell & $10^6$ & $\approx 4\cdot10^6$ & BFS: fine \\
values $1..10^6$, two operations & number & $10^6$ & $2\cdot10^6$ & BFS: fine \\
4-digit lock, turn one wheel $\pm1$ & 4 digits & $10^4$ & $8\cdot10^4$ & trivial \\
two tokens on an $8\times8$ board & 4-tuple & $4096$ & small & trivial \\
permutations of 8, adjacent swaps & permutation & $40\,320$ & $282\,240$ & fine (map or rank) \\
20 lamps, toggle one & bitmask & $1\,048\,576$ & $20\,971\,520$ & fine, $\approx 0.2$ s \\
all words of length 5 & string & $11\,881\,376$ & $\times 125$ & dictionary words only \\
values up to $10^{18}$ & number & $10^{18}$ & --- & impossible \\
\bottomrule
\end{tabularx}
\end{table}

If the estimate is far above $10^8$, the plain traversal is not the solution: look for a smaller state, a bound on the useful values, or mathematics (for the last row: a greedy argument, or Chapter~11's binary lifting when the graph is functional).

%------------------------------------------------------------
\subsubsection*{Recognising the shape from the statement}

The constraints and the wording often reveal \emph{which kind} of graph you have, and with it the algorithm family:

\begin{diagrambox}
\begin{tikzpicture}[
    clue/.style={rectangle, rounded corners=2pt, draw=accentorange, fill=memobg, font=\scriptsize, align=left, text width=5.6cm, minimum height=7mm},
    shp/.style={rectangle, rounded corners=2pt, draw=primaryblue, fill=theorybg, font=\scriptsize\bfseries, align=center, text width=2.4cm, minimum height=7mm},
    algo/.style={font=\scriptsize\color{darkgray}, align=left, text width=3.4cm, anchor=west}]
    \foreach \c/\s/\a [count=\i] in {
        {``$n$ vertices, $n-1$ edges, connected''}/tree/{unique paths; Ch.~13},
        {``$a_i$ is the next / the target of $i$''}/functional graph/{follow arrows; Pattern~5},
        {``prerequisites'', ``must come before''}/DAG/{topological order; Ch.~5},
        {``two teams'', ``enemies in different groups''}/bipartite/{2-colouring; Ch.~4},
        {a board of \texttt{.} and \texttt{\#}, ``islands''}/grid graph/{flood fill, BFS; Pattern~3},
        {``minimum number of moves / operations''}/state graph/{BFS on states; Pattern~4},
        {``how many friends'', ``trusted by everyone''}/degrees only/{counting; Pattern~1},
        {``groups'', ``connected?'', ``can reach''}/reachability/{DFS/BFS, DSU; Pattern~2}}{
        \node[clue] (c\i) at (0,-\i*0.95) {\c};
        \node[shp] (s\i) at (5.0,-\i*0.95) {\s};
        \node[algo] at (6.4,-\i*0.95) {\a};
        \draw[flow] (c\i) -- (s\i);
    }
    \node[font=\scriptsize\bfseries] at (0,-0.25) {clue in the statement};
    \node[font=\scriptsize\bfseries] at (5.0,-0.25) {shape};
    \node[font=\scriptsize\bfseries, anchor=west] at (6.4,-0.25) {first idea};
\end{tikzpicture}
\figcap{From wording to shape to algorithm. These are first guesses, not laws --- but they are right surprisingly often.}
\end{diagrambox}

\begin{diagrambox}
\begin{tikzpicture}[x=1cm, y=1cm]
    \begin{scope}
        \node[gv, minimum size=6mm] (a) at (0.9,0.9) {};
        \node[gv, minimum size=6mm] (b) at (0.2,0) {};
        \node[gv, minimum size=6mm] (c) at (1.6,0) {};
        \node[gv, minimum size=6mm] (d) at (1.6,-0.9) {};
        \draw[gtree] (a) -- (b); \draw[gtree] (a) -- (c); \draw[gtree] (c) -- (d);
        \node[font=\scriptsize\bfseries] at (0.9,-1.6) {tree};
    \end{scope}
    \begin{scope}[xshift=2.9cm]
        \node[gvsel, minimum size=6mm] (a) at (0.4,0.7) {};
        \node[gvsel, minimum size=6mm] (b) at (1.4,0.7) {};
        \node[gvsel, minimum size=6mm] (c) at (0.9,-0.1) {};
        \node[gv, minimum size=6mm] (d) at (0.9,-1.0) {};
        \draw[gde] (a) -- (b); \draw[gde] (b) -- (c); \draw[gde] (c) -- (a); \draw[gde] (d) -- (c);
        \node[font=\scriptsize\bfseries] at (0.9,-1.6) {functional};
    \end{scope}
    \begin{scope}[xshift=5.6cm]
        \node[gv, minimum size=6mm] (a) at (0,0.5) {};
        \node[gv, minimum size=6mm] (b) at (0,-0.6) {};
        \node[gv, minimum size=6mm] (c) at (1,0) {};
        \node[gv, minimum size=6mm] (d) at (2,0) {};
        \draw[gde] (a) -- (c); \draw[gde] (b) -- (c); \draw[gde] (c) -- (d); \draw[gde] (a) to[bend left=30] (d);
        \node[font=\scriptsize\bfseries] at (1,-1.6) {DAG};
    \end{scope}
    \begin{scope}[xshift=8.6cm]
        \node[gvA, minimum size=6mm] (a) at (0,0.6) {};
        \node[gvA, minimum size=6mm] (b) at (0,-0.6) {};
        \node[gvB, minimum size=6mm] (c) at (1.2,0.6) {};
        \node[gvB, minimum size=6mm] (d) at (1.2,-0.6) {};
        \draw[ge] (a) -- (c); \draw[ge] (a) -- (d); \draw[ge] (b) -- (c);
        \node[font=\scriptsize\bfseries] at (0.6,-1.6) {bipartite};
    \end{scope}
    \begin{scope}[xshift=11.2cm, yshift=0.5cm]
        \foreach \r in {0,1,2} \foreach \c in {0,1,2} \node[gcell, minimum size=5mm] at (\c*0.5,-\r*0.5) {};
        \node[gwall, minimum size=5mm] at (0.5,-0.5) {};
        \node[font=\scriptsize\bfseries] at (0.5,-2.1) {grid};
    \end{scope}
\end{tikzpicture}
\figcap{The shapes of this chapter's toolbox, in miniature (the purple vertices of the functional graph form its cycle). Learn to see them in a statement before you see them in a drawing.}
\end{diagrambox}

%------------------------------------------------------------
\subsubsection*{Worked example: modelling a statement}

\begin{examplebox}[A kingdom of towns]
\emph{``The kingdom has $n\le 2\cdot10^5$ towns and $n-1$ roads; every town can be reached from the capital (town 1). For each town, how many roads separate it from the capital?''}

\step{Vertex:} a town. \step{Edge:} a road, undirected, unweighted. \step{Shape:} $n-1$ edges and connected $\Rightarrow$ a \emph{tree}. \step{Question:} the number of roads from vertex 1 to every vertex $=$ BFS distances (or DFS depths: in a tree the path is unique, so DFS gives the shortest distance too). \step{Size:} $|V|=2\cdot10^5$, $|E|=2\cdot10^5-1$, so $\bigO{n}$ is instant; adjacency list, not a matrix; BFS or iterative DFS, because a path-shaped tree is $2\cdot10^5$ levels deep.
\end{examplebox}

%------------------------------------------------------------
\subsubsection*{Check your understanding}

For each situation give the vertex, the edge and an estimate of $|V|$.
\begin{enumerate}[leftmargin=*]
    \item A knight on a $100\times100$ board must reach a target square in as few moves as possible.
    \item A robot on an $R\times C$ grid ($R,C\le 50$) collects up to 5 different keys that open doors.
    \item ``$n\le 10^5$ students; student $i$ copies the homework from student $p_i$.''
    \item From $x$ you may go to $x+3$ or $x-5$, but must always stay inside $[0,10^6]$; can you get from $a$ to $b$?
    \item A statement has $n\le 10^5$ cities, and every pair of cities is joined by a road \emph{except} for $k$ listed pairs. Why can you not build the adjacency list?
\end{enumerate}

\begin{solutionbox}[Answers]
\begin{enumerate}[leftmargin=*]
    \item Vertex $=$ square, edge $=$ knight move (up to 8 per square, undirected), $|V|=10^4$; BFS.
    \item Vertex $=$ (row, column, set of keys held), the keys as a 5-bit mask: $|V|=50\cdot50\cdot32=80\,000$. The position alone is not enough: the same cell with different keys has a different future.
    \item Vertex $=$ student, arc $i\to p_i$: every vertex has out-degree 1 --- a functional graph (Pattern~5), $|V|=10^5$.
    \item Vertex $=$ integer in $[0,10^6]$, arcs $x\to x+3$ and $x\to x-5$ when they stay in range: $|V|=10^6+1$, at most $2\cdot10^6$ arcs. A BFS or DFS from $a$ answers it (directed reachability).
    \item The graph is almost complete: about $n^2/2 = 5\cdot10^9$ edges for $n=10^5$, far beyond memory. Work with the $k$ \emph{missing} pairs (the complement) and a cleverer traversal instead.
\end{enumerate}
\end{solutionbox}
